\section{Conclusões Parciais e Próximos Passos}
\label{sec:conclusoes}

Até este ponto, o projeto estabeleceu uma base reprodutível (M0) e validou a
qualidade dos dados (M1). Os principais resultados parciais são:
\begin{itemize}[noitemsep]
  \item Infraestrutura com configuração centralizada, \emph{seeds} fixas e proveniência
        documentada de cada hiperparâmetro, separando escolhas fundamentadas de valores
        provisórios.
  \item Quatro séries íntegras (3724 pregões, sem NaN), com anomalias reais já mapeadas
        como âncoras para a validação narrativa.
  \item Esclarecimento do caso AMER3, com decisão registrada de correção do artefato de
        grupamento na próxima etapa.
\end{itemize}

\paragraph{Próximos passos (M2 --- Pré-processamento).} Implementar o cálculo de
log-retornos, a correção do fator de grupamento do AMER3, o \emph{split} temporal com
ajuste do \texttt{MinMaxScaler} restrito ao treino e a geração das janelas deslizantes.
As etapas subsequentes cobrirão o treino do modelo (M3), os limiares de detecção (M4), a
avaliação quantitativa por injeção sintética (M5) e a correlação com eventos e a
comparação setorial (M6).

\paragraph{Registro de decisões.} Todas as decisões de configuração estão documentadas
como ADRs em \texttt{docs/adr/}. Este relatório é atualizado a cada fechamento de
\emph{milestone}, consolidando decisões, fontes, observações e conclusões para compor o
artigo final.
